\documentclass[12pt,a4paper]{letter}
\usepackage[utf8]{inputenc}
\usepackage[margin=1in]{geometry}
\usepackage{times}
\usepackage{setspace}
\onehalfspacing

% Letter details
\name{Deborah Akuoko}

\begin{document}

\begin{letter}{Professor Guangzhao Mao\\Head of the School of Engineering\\Chair Professor of Materials Engineering\\School of Engineering\\The University of Edinburgh\\King's Buildings\\Edinburgh\\EH9 3FB\\United Kingdom}

\date{\today}

\opening{Dear Professor Mao,}

I hope this letter finds you well. I am writing to request your consideration in providing a letter of recommendation to support my application for the UK Global Talent visa.

I have already submitted my application with a comprehensive recommendation letter from my academic supervisor, Dr. Istvan Gyongy of the University of Edinburgh, and with the full support of the University itself. However, UK Visas and Immigration has responded to my application by requesting an additional letter of recommendation from an eminent person in the UK. This official request, which I have received via email, is attached for your review along with Dr. Gyongy's recommendation letter.

The reason I am approaching you, Professor Mao, is precisely because UKVI specifically requires a recommendation from an 'eminent' person based in the UK. By definition, an eminent person is one who has achieved distinguished recognition and standing in their field, typically evidenced by significant honours, leadership positions in prestigious institutions, and substantial contributions to their discipline. Your distinguished career exemplifies these criteria: as Head of the School of Engineering and Chair Professor of Materials Engineering at the University of Edinburgh, you lead one of the most prestigious engineering schools in the United Kingdom. Your pioneering research in nanotechnology and materials engineering—particularly your groundbreaking work in electrocrystallization and nanosensors, as well as your innovative approaches to drug delivery systems targeting the central nervous system—has earned you international recognition. Your prestigious accolades, including election as a Fellow of the American Institute of Chemical Engineers (AIChE), Fulbright Senior Scholarship, visiting professorship at the Max Planck Institute of Colloids and Interfaces, and a National Science Foundation CAREER Award, demonstrate your standing as a leading innovator in your field. Your previous leadership as Head of the School of Chemical Engineering at the University of New South Wales, Sydney, where you successfully increased research income, established new research centres, and enhanced industry engagement, further exemplifies the calibre of eminent person that UKVI seeks to provide such endorsements. Your standing in the UK and international academic community, combined with your extensive contributions to nanotechnology, materials engineering, and engineering education, makes you uniquely qualified to provide the authoritative assessment that UKVI requires.

I am particularly moved to reach out to you because of our direct connection through the University of Edinburgh. As Head of the School of Engineering, you lead the institution where I have pursued my doctoral research. As a current doctoral candidate at Edinburgh, I have had the privilege of advancing my research under the auspices of this esteemed institution, and I believe that your understanding of the School's values, research excellence, and commitment to innovation would enable you to appreciate the significance of the work I have undertaken here. Having recently submitted my PhD thesis at the University of Edinburgh, my research has resulted in both patent applications and publication drafts that demonstrate my contributions to the field. The intellectual property and scholarly work emerging from my doctoral research reflects my competence and expertise, as well as my potential to make meaningful contributions to the academic and research landscape in the United Kingdom, particularly in areas that align with the School of Engineering's research themes.

I must, with some trepidation, share that my current circumstances are marked by considerable urgency. At present, I have no source of income, and the successful outcome of this visa application is not merely a professional aspiration but a matter of personal survival. The prompt reception of the visa would go a long way to grant me financial opportunities that would enable me to help and support an ailing close relative who depends on my assistance. 

I must also acknowledge, with candour, that this matter has become more urgent for me than the completion of my PhD, for I cannot take pride in earning a prestigious degree when I find myself incapacitated as a result of having no right to work. The opportunity to continue my research and contribute to the UK's academic community through the Global Talent visa would provide not only the stability I so desperately need but also the platform to realise the potential that my supervisors and institution have recognised in my work.

I understand that requests of this nature require careful consideration, and I would be profoundly grateful for any support you might be able to offer. Your endorsement would carry significant weight in the visa application process and would represent a transformative opportunity for my career and personal circumstances.

I have enclosed the UKVI email request, the recommendation letter from Dr. Istvan Gyongy, along with a single document briefing my working papers and patents, should you wish to review the evidence of my work and the support I have already received from my academic institution.

Thank you most sincerely for your time and consideration. I would be delighted to provide any additional information you might require and would welcome the opportunity to discuss my research and aspirations with you at your convenience.

\closing{Yours sincerely,}

\signature{Deborah Akuoko}

\encl{UKVI email request for additional recommendation\\Recommendation letter from Dr. Istvan Gyongy\\Document briefing working papers and patents}

\end{letter}

\end{document}

